%%%%%%%%%%%%%%%%%
% This is an example CV created using altacv.cls (v1.3, 10 May 2020) written by
% LianTze Lim (liantze@gmail.com), based on the
% Cv created by BusinessInsider at http://www.businessinsider.my/a-sample-resume-for-marissa-mayer-2016-7/?r=US&IR=T
%
%% It may be distributed and/or modified under the
%% conditions of the LaTeX Project Public License, either version 1.3
%% of this license or (at your option) any later version.
%% The latest version of this license is in
%%    http://www.latex-project.org/lppl.txt
%% and version 1.3 or later is part of all distributions of LaTeX
%% version 2003/12/01 or later.
%%%%%%%%%%%%%%%%

%% If you are using \orcid or academicons
%% icons, make sure you have the academicons
%% option here, and compile with XeLaTeX
%% or LuaLaTeX.
% \documentclass[10pt,a4paper,academicons]{altacv}

%% Use the "normalphoto" option if you want a normal photo instead of cropped to a circle
% \documentclass[10pt,a4paper,normalphoto]{altacv}

\documentclass[10pt,a4paper,ragged2e,withhyper]{altacv}

%% AltaCV uses the fontawesome5 and academicon fonts
%% and packages.
%% See http://texdoc.net/pkg/fontawesome5 and http://texdoc.net/pkg/academicons for full list of symbols. You MUST compile with XeLaTeX or LuaLaTeX if you want to use academicons.

% Change the page layout if you need to
\geometry{left=0.9cm,right=0.9cm,top=0.9cm,bottom=0.9cm,columnsep=0.8cm}

% The paracol package lets you typeset columns of text in parallel
\usepackage{paracol}


% Change the font if you want to, depending on whether
% you're using pdflatex or xelatex/lualatex
\ifxetexorluatex
  % If using xelatex or lualatex:
  \setmainfont{Lato}
\else
  % If using pdflatex:
  \usepackage[default]{lato}
\fi

% Change the colours if you want to
\definecolor{arsenic}{rgb}{0.23, 0.27, 0.29}
\definecolor{SlateGrey}{HTML}{2E2E2E}
\definecolor{LightGrey}{HTML}{666666}
\definecolor{Blue}{HTML}{1EA5E8}
\colorlet{heading}{arsenic}
\colorlet{accent}{arsenic}
\colorlet{emphasis}{SlateGrey}
\colorlet{body}{LightGrey}
\colorlet{subheading}{Blue}

% Change some fonts, if necessary
% \renewcommand{\namefont}{\Huge\rmfamily\bfseries}
% \renewcommand{\personalinfofont}{\footnotesize}
% \renewcommand{\cvsectionfont}{\LARGE\rmfamily\bfseries}
% \renewcommand{\cvsubsectionfont}{\large\bfseries}

% Change the bullets for itemize and rating marker
% for \cvskill if you want to
\renewcommand{\itemmarker}{{\small\textbullet}}
\renewcommand{\ratingmarker}{\faCircle}
\newcommand*\googlescholaricon{\raisebox{-0.3em}{\includegraphics[width=1em]{google-scholar.pdf}}}

%% sample.bib contains your publications
\addbibresource{sample.bib}

\begin{document}
\name{Shreyas Kalvankar}
\tagline{}
%% You can add multiple photos on the left or right
%\photoR{2.5cm}{mmayer-wikipedia-cc-by-2_0}
% \photoL{2cm}{Yacht_High,Suitcase_High}
\personalinfo{%
  % Not all of these are required!
  % You can add your own with \printinfo{symbol}{detail}
  % \printinfo{\faGlobe}{obi-wan-shinobi.github.io}
  \homepage{obi-wan-shinobi.github.io}
  \email{shreyaskalvankar@gmail.com}
  \phone{+919423555723}
  \location{Maharashtra, India}
  \linkedin{linkedin.com/in/shreyas-kalvankar}
  \github{github.com/obi-wan-shinobi}
  \printinfo{\googlescholaricon}{Scholar Profile \href{https://scholar.google.com/citations?user=W8-T__UAAAAJ&hl=en}{[link]}}
%   \github{github.com/mmayer} % I'm just making this up though.
%   \orcid{orcid.org/0000-0000-0000-0000} % Obviously making this up too. If you want to use this field (and also other academicons symbols), add "academicons" option to \documentclass{altacv}
  %% You MUST add the academicons option to \documentclass, then compile with LuaLaTeX or XeLaTeX, if you want to use \orcid or other academicons commands.
  % \orcid{0000-0000-0000-0000}
  %% You can add your own arbtrary detail with
  %% \printinfo{symbol}{detail}[optional hyperlink prefix]
  % \printinfo{\faPaw}{Hey ho!}
  %% Or you can declare your own field with
  %% \NewInfoFiled{fieldname}{symbol}[optional hyperlink prefix] and use it:
  % \NewInfoField{gitlab}{\faGitlab}[https://gitlab.com/]
  % \gitlab{your_id}
}

\makecvheader

%% Depending on your tastes, you may want to make fonts of itemize environments slightly smaller
\AtBeginEnvironment{itemize}{\small}

%% Set the left/right column width ratio to 6:4.
\columnratio{0.6}

% Start a 2-column paracol. Both the left and right columns will automatically
% break across pages if things get too long.
\begin{paracol}{2}

\cvsection{Experience}
\cvevent{Software Developer}{Dalton Maag Ltd.}{November 2021 - Present}{London, United Kingdom}
\begin{itemize}
\begin{footnotesize}
    \item Developed a Python-based Proof of Concept (POC) utilizing Genetic Algorithms to \textbf{automate the generation of thousands of CJK font glyphs}, with the projected potential to significantly reduce production time from \textbf{months} to \textbf{minutes}.
    % \item Developed PriceBot, a web application with Ruby on Rails back-end and a VueJS \& Typescript front-end to simplify typeface pricing.
    \item Accelerated quote generation from \textbf{hours} to \textbf{seconds} by automating pricing with PriceBot, refining pricing models in Typescript and Ruby on Rails.
    % \item Implemented and fine-tuned pricing models in Typescript, allowing PriceBot to automate quote generation within \textbf{seconds}, a drastic improvement over the previous manual process that typically took \textbf{hours}.
    \item Implemented an application-health-tracking feature in Ruby on Rails to monitor pricing model states, issuing warnings to end users when pricing diverged significantly from expected figures, thereby preventing revenue loss by saving \textbf{hundreds of thousands of GBP} in response to the challenge of potential inaccuracies in project quotes.
    \item Containerized the application-health-tracking feature and seamlessly integrated it with GitLab CI/CD for continuous monitoring, enhancing operational efficiency.
    \item Developed glyph data models for Arabic, Greek, Cyrillic, and Devanagari, ensuring precise pricing for non-Latin projects, contributing to an increased client base.
    \item Implemented a streamlined process in Typescript utilizing graphs theory to efficiently generate project plans, \textbf{reducing planning time from hours to a few minutes}, and delivering precise turnaround times to clients.
    \item Contributed to the development of DSedit, an ElectronJS application with a VueJS frontend, serving as a unified platform housing all internal tools and scripts, aimed at providing seamless accessibility and enhancing operational efficiency.
\end{footnotesize}
\end{itemize}
\vspace{5px}
% \divider
\cvevent{Machine Learning Research Scientist (Consultant)}{Relfor Labs Pvt. Ltd.}{September 2022 - Present}{Pune, India}
\begin{itemize}
\begin{footnotesize}
    \item Set up the ML training pipeline using PyTorch lightning on Nvidia DGX A100, with automatic hyperparameter tuning using Optuna, expediting experimentation, making it \textbf{$\sim$50\%} faster.
    \item Revamped ETL pipeline to optimize CPU and storage usage, resulting in an \textbf{8x} speed improvement and a \textbf{30\%} boost in storage efficiency while processing terabytes of data.  
\end{footnotesize}
\end{itemize}
\cvevent{Machine Learning Engineer}{Relfor Labs Pvt. Ltd.}{August 2021 - November 2021}{Pune, India}
\begin{itemize}
\begin{footnotesize}
    \item Designed multiple novel deep convolutional neural network architectures for audio data classification in PyTorch, which beat state-of-the-art models with \textbf{~98.6\%} accuracy and \textbf{$\sim$0.98} F1-score.
   \item Implemented various threshold optimization techniques in SKLearn and PyTorch, achieving a precision of \textbf{$\sim$98\%} while maintaining high accuracy \textbf{>98\%}.
\end{footnotesize}
\end{itemize}
% \divider
\vspace{5px}
\cvevent{Software Development Intern}{FinIQ Consulting India Pvt. Ltd.}{May 2020 - June 2020}{Nashik, India}
\begin{itemize}
\begin{footnotesize}
\item Set up an online platform for Forex trading using AngularJS as a new feature for the customers.
\item Created an open-source Python module for stress testing CPU and memory with variable load which was later integrated into the company's cloud platforms' testing pipeline.
\end{footnotesize}

\end{itemize}

\begin{comment}
\cvsection{Technical Skills}
\cvskill{C/C++, Python, Java}{4}
\cvskill{Deep Learning}{4}
\cvskill{Computer Vision}{4}
\cvskill{Machine Learning}{3}
\cvskill{Databases}{3}
\cvskill{Robotics}{3}
\cvskill{Web Development}{2}
\end{comment}


\cvsection{Technical skills}
% Python, C, C++, Javascript(Babel ES6, Typescript, Webpack, Mocha, Chai),
% ReactJs, VueJS, HTML, CSS, Ruby on Rails,
% MongoDB, PostgreSQL, Docker, Keras, Tensorflow, PyTorch
\begin{itemize}
\begin{footnotesize}
    \item \textbf{Computer Languages} : C, C++, Python, Java (J2EE)
    \item \textbf{Web Development} : Javascript (Babel ES6, Typescript, Webpack, Mocha, Chai), HTML, CSS, Ruby on Rails, ReactJS, VueJS
    \item \textbf{ML Frameworks} : Keras, Tensorflow, PyTorch
    \item \textbf{Database \& Big Data} : mongoDB, Postgres, Spark
    \item \textbf{Containerization} : Docker
    \item \textbf{Cloud}: Google Cloud Platform, Cloud Queue, GCP Buckets, Google Cloud Functions
    % \item \textbf{Machine Learning Frameworks} : Octave, Sci-kit Learn
    % \item \textbf{Embedded Software Programming} : Arduino, Raspberry Pi, Teensy
    % \item \textbf{Version Control} : Git, GitHub
\end{footnotesize}
\end{itemize}
% \cvsection{Positions of Responsibility}
% \cvevent{Software Developer}{Team Vector, ABU Robocon 2019}{August 2018 - April 2019}{}
% \begin{itemize}
%     \item Assigned to build and code a quadruped robot and a wheeled robot with dynamic locomotive abilities for ABU Robocon 2019
% \end{itemize}

% \divider

% \cvevent{Mentor}{Team Vector, ABU Robocon 2020}{August 2019 - April 2020}{}
% \begin{itemize}
%     \item Helped and guided junior members of the team in building robots that could efficiently handle locomotion and throwing, catching and kicking a football
% \end{itemize}

\begin{comment}
\cvsection{Strengths}
\cvtag{Deep Learning}
\cvtag{Machine Learning \& Data Science}\\
\cvtag{Software Development}
\cvtag{Researcher}
\cvtag{Hardworking}\\
\cvtag{Adaptable}
\cvtag{Communication skills}
\end{comment}
%% Switch to the right column. This will now automatically move to the second
%% page if the content is too long.
\switchcolumn

\cvsection{Education}
\cvevent{B.E (Computer Engineering)}{K.K. Wagh Institute of Engineering Education and Research}{2017-2021}{Nashik}
\begin{itemize}
    \item CGPA: 9.7/10 (Rank 1)
\end{itemize}

\cvsection{Personal}

    \cvsubsection{THE GALAXY ZOO PROJECT}
    \begin{itemize}
        \item Developed a CNN in Tensorflow for vote fraction predictions of 37 galaxy features from the Galaxy Zoo decision tree with an rmse score of \textbf{0.07765}, ranking us in the \textbf{top 3} on the public leaderboard
        \item Also developed a CNN for classification of galaxies into 7 classes based on their morphologies with an accuracy of \textbf{93.7\%} and an F1 score of \textbf{0.8857}
    \end{itemize}
\begin{comment}
\cvsubsection{THE SCHRODPY PROJECT}
    \begin{itemize}
        \item The project is in early phases, which plans at providing an efficient simulation of quantum systems
        \item Applied various methods to numerically solve Schrödinger's Time Independent wave equation for finding stationary states of the particle in an Infinite Potential Well
        \item Prospects: An efficient and community python package for simulating quantum many-body systems using tensor networks
    \end{itemize}
\end{comment}
\cvsubsection{THE EINSTEINPY PROJECT}
\begin{itemize}
    \item An open source community python package for general relativity
    \item \textbf{Contributions}:
    \begin{itemize}
        \item Added Reissner–Nordström metric: a static solution to the Einstein-Maxwell field equations, into the code %(\textit{PR: \#462 Issue: \#309})
        \item Corrections in the Kerr-Newman and Kerr metrics classes
        \item Added calculations of event horizon and ergosphere for a Kerr-Newman blackhole %(\textit{PR: \#472 Issue: \#109}) 
        \item \href{https://doi.org/10.5281/zenodo.4445219}{DOI: 10.5281/zenodo.4445219}
    \end{itemize}
\end{itemize}
\cvsection{Publications}
\nocite{*}
\printbibliography[heading=pubtype,title={\printinfo{\faFile*[regular]}{Journal Articles}}, type=misc]


\end{paracol}

\end{document}
